\subtituloA{Sintaxis de apertura del script}

\definicion{Atributo type}{El atributo \texttt{type="text/javascript"} en la etiqueta de apertura informa al navegador que el contenido del script está formateado como texto javascript, no como html.}

\begin{codigo}
 
\cmd{<script type="text/javascript">} 
// Código javascript va aquí 
\cmd{</script>}
\bend{codigo}

\nota{Para el examen}{
¿Cuál es la diferencia práctica entre colocar un script en \texttt{<head>} versus en \texttt{<body>} en términos de disponibilidad de elementos del documento?

\textbf{Respuesta:} En \texttt{<head>} el script se ejecuta antes de que el cuerpo de la página esté cargado, por lo que no puede manipular elementos del cuerpo; en \texttt{<body>} los elementos ya existen y son accesibles por el script.
}}
\end{codigo}


